\documentclass[11pt,a4paper]{article}
\usepackage[T1]{fontenc}
\usepackage[utf8]{inputenc}
\usepackage{amsmath,amssymb,amsthm}
\usepackage[margin=1in]{geometry}
\usepackage[colorlinks=true,linkcolor=blue,urlcolor=blue]{hyperref}
\theoremstyle{plain}
\newtheorem{theorem}{Theorem}
\newtheorem{lemma}[theorem]{Lemma}
\newtheorem{proposition}[theorem]{Proposition}
\newtheorem{corollary}[theorem]{Corollary}
\theoremstyle{definition}
\newtheorem{remark}[theorem]{Remark}

\title{The empirical recurrence for OEIS A282796, proved from an exact model}
\author{Adrian Perez Fontelles\\ \small Independent researcher}
\date{6 September 2026}
\begin{document}
\maketitle

\begin{abstract}
OEIS A282796 carries an empirical recurrence of order $4$ contributed by Colin Barker. It is true, and it is
decidable rather than empirical. The entry's sequence is given exactly by a certified growth pattern of the automaton's axis word, from
which $a$ provably satisfies a MONIC linear recurrence of order $S=5$, derived below and not
assumed. The residual of the conjectured recurrence satisfies the same one, so whether it
annihilates $a$ is settled by a finite exact computation. No transfer matrix is involved and
the count is not a walk.
\end{abstract}

\noindent\small 2020 Mathematics Subject Classification. 68Q80, 11B85, 05A15.\normalsize

\section{The conjecture}

OEIS A282796 is ``Binary representation of the x-axis, from the left edge to the origin, of the n-th stage of growth of the two-dimensional cellular automaton defined by "Rule 505", based on the 5-celled von Neumann neighborhood.''. It has offset $0$ and begins
\[
1, 0, 111, 100, 11111, 10100, 1111111, 1010100,\ \dots
\]
The entry states, as a conjecture contributed by its author and never marked settled:
\begin{quote}
Conjectures from \_Colin Barker\_, Feb 21 2017: (Start) a(n) = 101*a(n-2) - 100*a(n-4) for n>3.
\end{quote}
As of the ``Last modified'' line on the live entry (Nov 05 2025 15:22:37, revision 11)
this is still recorded as empirical, and nothing on the entry records it as proved.

\section{The value is a certified growth}

The entry reads an axis or a diagonal of the automaton at stage $n$ as a string of digits and
takes its numeric value. Nothing is being counted, so there is no digraph. What settles the
sequence is that the CONFIGURATION is periodic in time:
\[
C(n+0)\;=\;C(n)\qquad\text{as configurations of the whole plane, for every } n\ge 0.
\]
This was checked directly, and one instance of it is all that is needed: the automaton is a
deterministic map on configurations, so $C(n_0+0)=C(n_0)$ forces $C(n_0+k\cdot0)=C(n_0)$ for
every $k$, and likewise in each residue class. No locality argument and no induction over
boundary windows is involved.

Getting the background right is the whole of the check. A rule taking an empty von Neumann
neighbourhood to a live cell flips the ENTIRE plane, and rule 0 is read
in that light: comparing two stages by padding the smaller with zeros would compare two
different things.

The reading follows at once. The axis at stage $n$ is the window $|x|\le n$ of a configuration
that is not changing, so the window widening by one cell on each side per step gives
\[
w(n+0)\;=\;L_r + w(n) + R_r
\]
with $L_r$ and $R_r$ fixed words depending only on the residue $r$ of $n$ modulo $0$ --- an
identity that is now a consequence rather than an observation. Concatenation is multiplication
by a power of the base $B$ and addition, so along a class
\[
a(n+0)\;=\;v(L_r)\,B^{|w(n)|+|R_r|} + a(n)\,B^{|R_r|} + v(R_r).
\]

\section{The bound}

In the shift $T=S^{\,0}$ that is a first-order recurrence with a geometric forcing term and a
constant, killed by $(T-B^{|R_r|})(T-B^{|L_r|+|R_r|})(T-1)$. Taking the DISTINCT roots that
occur across the residue classes, pulling each back to $S^{\,0}-\lambda$, and allowing the
pre-period to raise the numerator's degree gives a monic annihilator of order $S=5$.

\section{The proof}

\begin{lemma}\label{lem:crit}
Let $A(z)=z^S-\sum_{j<S}\alpha_jz^{j}$ be the monic annihilator derived above and
$q(z)=z^{4}-\sum_ic_iz^{4-i}$ the characteristic polynomial of the conjectured
recurrence. The residual $r(n)=a(n)-\sum_ic_ia(n-i)$ is a linear combination of shifts of $a$,
so $A$ annihilates $r$ as well. Since $A(0)\ne0$ the recurrence it gives runs in both
directions, and $S$ consecutive zeros of $r$ force $r$ to vanish at every later index.
\end{lemma}

\begin{proof}
$A$ annihilates $a$, hence every shift of $a$, hence any linear combination of shifts; that is
$r$. A monic recurrence with nonzero constant term determines each term from the $S$ before it
and each from the $S$ after it, so a run of $S$ zeros propagates.
\end{proof}

\begin{theorem}
\[
a(n)\;=\;101\,a(n-2) - 100\,a(n-4)
\]
for every $n>0$.
\end{theorem}

\begin{proof}
The terms $a(0),\dots$ were computed exactly in integer arithmetic from the model, extended by
$A$ where the model itself was not evaluated, and $r(n)$ formed for each index. Those integers
vanish from $n=0+1$ onwards, and the run exceeds $S+4$, which by
Lemma~\ref{lem:crit} settles every larger index at once.
\end{proof}

\section{Verification}

Three checks, each independent of the algebra above.

First, the model reproduces all $20$ terms the entry publishes, exactly, in integer
arithmetic. This is what ties the model to the entry: it is built by reading the entry's
English, and a misreading gives different counts.

Second, the conjectured recurrence was evaluated directly on those published terms, with no
model involved, and holds wherever the proved range and the data overlap.

Third, the derived annihilator $A$ was itself tested on terms it had not been given: the model
was evaluated beyond the $S$ values that determine it and $A$ reproduced them. A bound that is
too small fails this test, which is why it is run.

\begin{thebibliography}{9}
\bibitem{oeis} The OEIS Foundation, \emph{The On-Line Encyclopedia of Integer
Sequences}, \url{https://oeis.org}, 2026. Sequence A282796.
\bibitem{beckrobins} M.~Beck and S.~Robins, \emph{Computing the Continuous Discretely},
2nd ed., Springer, 2015. (Ehrhart quasi-polynomials and their periods.)
\bibitem{stanley} R.~P.~Stanley, \emph{Enumerative Combinatorics}, Volume 1, 2nd ed.,
Cambridge University Press, 2011.
\bibitem{ds} F.~Diamond and J.~Shurman, \emph{A First Course in Modular Forms},
Springer GTM 228, 2005.
\end{thebibliography}

\end{document}
